% GENERATED FILE. Edit canonical lessons, then run npm run generate:books.
% canonical-source-hash: fnv1a64:fb8e509823414a5a
% canonical-lessons: HI-C339-rajkumar, HI-C339-bhes, HI-C339-parakhna, HI-C339-chal, HI-C339-lubhana

\chapter{The Prince in Disguise}
\label{ch:hi-the-prince-in-disguise}
\hlchaptermodality{339}

\begin{chapteropening}
\textbf{By the end of this chapter:} \emph{I can tell how a prince tested the people.}

Complete the last lesson of chapter 339: I can tell how a prince tested the people.
\end{chapteropening}

\section[\texorpdfstring{\hi{राजकुमार} (rājkumār)}{राजकुमार (rājkumār)}]{\hi{राजकुमार} (rājkumār) — a prince}

\label{lesson:HI-C339-rajkumar}



\begin{quote}
Before the new one: say the Hindi for a witch, then the Hindi for a miracle.
\end{quote}

\subsection*{You'll want to know: \hi{राजकुमार}}
\textbf{\hi{राजकुमार}} — \emph{rājkumār} — \textquotedblleft{}a prince\textquotedblright{}.

\begin{quote}
\hi{राजकुमार} \hi{घोड़े} \hi{पर} \hi{जंगल} \hi{में} \hi{गया।}
\end{quote}

\begin{quote}
\emph{rājkumār ghoṛe par jaṅgal meṁ gayā.}
\end{quote}

\begin{quote}
The prince rode into the forest.
\end{quote}

Masculine.

\begin{cousinweb}
\textbf{\hi{राजकुमार}} (\emph{rājkumār}) is Sanskrit \emph{rājakumāra} (Monier-Williams), \textbf{\hi{राजा}}, a king, and \emph{kumāra}, a boy: a king's son.
\end{cousinweb}

\begin{grammarlens}[title={What you've built}]
One more word for telling what happened.
\end{grammarlens}

\subsection*{Your turn}
\begin{itemize}
\raggedright
  \item \emph{Say it:} \emph{rājkumār}
  \item \emph{Say it:} \emph{rājkumār}, once more
  \item \emph{Say it:} \emph{rājkumār ghoṛe par jaṅgal meṁ gayā.} — The prince rode into the forest.
  \item \emph{Recall:} say the Hindi for a witch, then the Hindi for a miracle, then say \emph{rājkumār} again
\end{itemize}

\begin{tcolorbox}[breakable,colback=teal!4,colframe=teal!35!black,title={Before you move on}]
What does \hi{राजकुमार} mean? (\textquotedblleft{}A prince\textquotedblright{}.) Say it once more.
\end{tcolorbox}
\section[\texorpdfstring{\hi{भेस} (bhes)}{भेस (bhes)}]{\hi{भेस} (bhes) — a disguise; a guise}

\label{lesson:HI-C339-bhes}



\begin{quote}
Before the new one: say the Hindi for a miracle, then the Hindi for a prince.
\end{quote}

\subsection*{You'll want to know: \hi{भेस}}
\textbf{\hi{भेस}} — \emph{bhes} — \textquotedblleft{}a disguise; a guise\textquotedblright{}.

\begin{quote}
\hi{राजकुमार} \hi{ने} \hi{किसान} \hi{का} \hi{भेस} \hi{बनाया।}
\end{quote}

\begin{quote}
\emph{rājkumār ne kisān kā bhes banāyā.}
\end{quote}

\begin{quote}
The prince disguised himself as a farmer.
\end{quote}

Masculine.

\begin{cousinweb}
\textbf{\hi{भेस}} (\emph{bhes}) goes back to Sanskrit \emph{veṣa}, \textquotedblleft{}dress, appearance\textquotedblright{} (CDIAL 12129), where CDIAL lists Hindi \emph{bhes}. \textbf{\hi{भेस} \hi{बदलना}} is to go in disguise.
\end{cousinweb}

\begin{grammarlens}[title={What you've built}]
One more word for telling what happened.
\end{grammarlens}

\subsection*{Your turn}
\begin{itemize}
\raggedright
  \item \emph{Say it:} \emph{bhes}
  \item \emph{Say it:} \emph{bhes}, once more
  \item \emph{Say it:} \emph{rājkumār ne kisān kā bhes banāyā.} — The prince disguised himself as a farmer.
  \item \emph{Recall:} say the Hindi for a miracle, then the Hindi for a prince, then say \emph{bhes} again
\end{itemize}

\begin{tcolorbox}[breakable,colback=teal!4,colframe=teal!35!black,title={Before you move on}]
What does \hi{भेस} mean? (\textquotedblleft{}A disguise; a guise\textquotedblright{}.) Say it once more.
\end{tcolorbox}
\section[\texorpdfstring{\hi{परखना} (parakhnā)}{परखना (parakhnā)}]{\hi{परखना} (parakhnā) — to test; to put to the test}

\label{lesson:HI-C339-parakhna}



\begin{quote}
Before the new one: say the Hindi for a prince, then the Hindi for a disguise.
\end{quote}

\subsection*{You'll want to know: \hi{परखना}}
\textbf{\hi{परखना}} — \emph{parakhnā} — \textquotedblleft{}to test; to put to the test\textquotedblright{}.

\begin{quote}
\hi{राजकुमार} \hi{ने} \hi{भेस} \hi{बदलकर} \hi{लोगों} \hi{को} \hi{परखा।}
\end{quote}

\begin{quote}
\emph{rājkumār ne bhes badalkar logoṁ ko parkhā.}
\end{quote}

\begin{quote}
The prince went in disguise and put the people to the test.
\end{quote}

Past, with \hi{ने} (he or she …): \hi{उसने} \hi{सोना} \hi{परखा}, \emph{usne sonā parkhā}.

\begin{cousinweb}
goes back to Sanskrit \emph{parīkṣate}, \textquotedblleft{}examines\textquotedblright{} (CDIAL 7904), the family of \textbf{\hi{परीक्षा}}, an examination. You \textbf{\hi{परखना}} gold, a coin or a friend.
\end{cousinweb}

\begin{grammarlens}[title={What you've built}]
One more word for telling what happened.
\end{grammarlens}

\subsection*{Your turn}
\begin{itemize}
\raggedright
  \item \emph{Say it:} \emph{parakhnā}
  \item \emph{Say it:} \emph{parakhnā}, once more
  \item \emph{Say it:} \emph{rājkumār ne bhes badalkar logoṁ ko parkhā.} — The prince went in disguise and put the people to the test.
  \item \emph{Recall:} say the Hindi for a prince, then the Hindi for a disguise, then say \emph{parakhnā} again
\end{itemize}

\begin{tcolorbox}[breakable,colback=teal!4,colframe=teal!35!black,title={Before you move on}]
What does \hi{परखना} mean? (\textquotedblleft{}To test; to put to the test\textquotedblright{}.) Say it once more.
\end{tcolorbox}
\section[\texorpdfstring{\hi{छल} (chal)}{छल (chal)}]{\hi{छल} (chal) — trickery; deceit}

\label{lesson:HI-C339-chal}



\begin{quote}
Before the new one: say the Hindi for a disguise, then the Hindi for to test.
\end{quote}

\subsection*{You'll want to know: \hi{छल}}
\textbf{\hi{छल}} — \emph{chal} — \textquotedblleft{}trickery; deceit\textquotedblright{}.

\begin{quote}
\hi{राक्षस} \hi{ने} \hi{छल} \hi{से} \hi{राजकुमार} \hi{को} \hi{पकड़ा।}
\end{quote}

\begin{quote}
\emph{rākṣas ne chal se rājkumār ko pakṛā.}
\end{quote}

\begin{quote}
The demon caught the prince by trickery.
\end{quote}

Masculine.

\begin{cousinweb}
\textbf{\hi{छल}} (\emph{chal}) is Sanskrit \emph{chala}, \textquotedblleft{}fraud, a trick\textquotedblright{} (Monier-Williams). It belongs to tales and to writing: \textbf{\hi{छल} \hi{से}}, by a trick.
\end{cousinweb}

\begin{grammarlens}[title={What you've built}]
One more word for telling what happened.
\end{grammarlens}

\subsection*{Your turn}
\begin{itemize}
\raggedright
  \item \emph{Say it:} \emph{chal}
  \item \emph{Say it:} \emph{chal}, once more
  \item \emph{Say it:} \emph{rākṣas ne chal se rājkumār ko pakṛā.} — The demon caught the prince by trickery.
  \item \emph{Recall:} say the Hindi for a disguise, then the Hindi for to test, then say \emph{chal} again
\end{itemize}

\begin{tcolorbox}[breakable,colback=teal!4,colframe=teal!35!black,title={Before you move on}]
What does \hi{छल} mean? (\textquotedblleft{}Trickery; deceit\textquotedblright{}.) Say it once more.
\end{tcolorbox}
\section[\texorpdfstring{\hi{लुभाना} (lubhānā)}{लुभाना (lubhānā)}]{\hi{लुभाना} (lubhānā) — to lure; to tempt}

\label{lesson:HI-C339-lubhana}



\begin{quote}
Before the new one: say the Hindi for to test, then the Hindi for trickery.
\end{quote}

\subsection*{You'll want to know: \hi{लुभाना}}
\textbf{\hi{लुभाना}} — \emph{lubhānā} — \textquotedblleft{}to lure; to tempt\textquotedblright{}.

\begin{quote}
\hi{डायन} \hi{ने} \hi{मिठाई} \hi{से} \hi{बच्चों} \hi{को} \hi{लुभाया।}
\end{quote}

\begin{quote}
\emph{ḍāyan ne miṭhāī se baccoṁ ko lubhāyā.}
\end{quote}

\begin{quote}
The witch lured the children with sweets.
\end{quote}

Past, with \hi{ने} (he or she …): \hi{उसने} \hi{बच्चे} \hi{को} \hi{लुभाया}, \emph{usne bacce ko lubhāyā}.

\begin{cousinweb}
\textbf{\hi{लुभाना}} (\emph{lubhānā}) goes back to Sanskrit \emph{lubhyati}, \textquotedblleft{}is eager, longs for\textquotedblright{} (CDIAL 11086): to make someone want something so much that they follow.
\end{cousinweb}

\begin{grammarlens}[title={What you've built}]
One more word for telling what happened.
\end{grammarlens}

\subsection*{Your turn}
\begin{itemize}
\raggedright
  \item \emph{Say it:} \emph{lubhānā}
  \item \emph{Say it:} \emph{lubhānā}, once more
  \item \emph{Say it:} \emph{ḍāyan ne miṭhāī se baccoṁ ko lubhāyā.} — The witch lured the children with sweets.
  \item \emph{Recall:} say the Hindi for to test, then the Hindi for trickery, then say \emph{lubhānā} again
\end{itemize}

\begin{tcolorbox}[breakable,colback=teal!4,colframe=teal!35!black,title={Before you move on}]
What does \hi{लुभाना} mean? (\textquotedblleft{}To lure; to tempt\textquotedblright{}.) Say it once more.
\end{tcolorbox}
